\documentclass[12pt]{article}
\usepackage{amsmath,amssymb}

\begin{document}
\section*{{2017 AMC 12B Problems/Problem 24}}
Source: AoPS Wiki

\subsection*{Solution}
Let $CD=1$ , $BC=x$ , and $AB=x^2$ . Note that $AB/BC=x$ . By the Pythagorean Theorem, $BD=\sqrt{x^2+1}$ . Since $\triangle BCD \sim \triangle ABC \sim \triangle CEB$ , the ratios of side lengths must be equal. Since $BC=x$ , $CE=\frac{x^2}{\sqrt{x^2+1}}$ and $BE=\frac{x}{\sqrt{x^2+1}}$ . Let F be a point on $\overline{BC}$ such that $\overline{EF}$ is an altitude of triangle $CEB$ . Note that $\triangle CEB \sim \triangle CFE \sim \triangle EFB$ . Therefore, $EF=\frac{x}{x^2+1}$ and $EG=\frac{x^3}{x^2+1}$ . Since $\overline{CF}$ and $\overline{BF}$ form altitudes of triangles $CED$ and $BEA$ , respectively, the areas of these triangles can be calculated. Additionally, the area of triangle $BEC$ can be calculated, as it is a right triangle. Solving for each of these yields: \begin{align*}
[BEC] &=[CED]=[BEA]=\frac{x^3}{2(x^2+1)} \\
[ABCD] &=[AED]+[DEC]+[CEB]+[BEA] \\
\frac{(BC)(AB+CD)}{2} &= 17*[CEB]+ [CEB] + [CEB] + [CEB] \\
\frac{x^3+x}{2} &= \frac{20x^3}{2(x^2+1)} \\
\frac{x}{x^2+1} &= \frac{20x^3}{x^2+1} \\
(x^2+1)^2 &=20x^2 \\
x^4-18x^2+1 &=0 \implies x^2=9+4\sqrt{5}=4+2(2\sqrt{5})+5 \\
\end{align*} Therefore, the answer is $\boxed{\textbf{(D) } 2+\sqrt{5}}$

\end{document}
